\section{Reference architecture of a self-aware steel structure}\label{sec:architecture}

The transition from a passively monitored steel frame to a self-aware structure requires more than additional sensors; it requires an organizing architecture in which raw physical measurements are progressively transformed into interpretable knowledge about the structural state. We propose a five-layer reference architecture (sensing, data acquisition and signal conditioning, physics-based digital twin, cognition, and decision and reporting), summarized in Fig.~\ref{fig:arch}, that maps directly onto the statistical pattern recognition (SPR) paradigm for SHM \citep{farrar2013shm} while extending it with a continuously updated digital twin and an embedded cognition layer. Each layer consumes the output of the one below it and returns feedback to the layers beneath, so that the architecture operates as a closed perceive--process--interpret--report loop rather than a one-way data pipeline. The physics-based and data-driven components are not alternatives but co-operating elements of a hybrid, physics-informed system \citep{brunton2019data}, in which law-based models supply features and constraints to the machine-learning layer, and the machine-learning layer in turn supplies updated parameters and residuals back to the physics model \citep{chiachio2022structural,zou2026deep}.

\subsection{Sensing layer}\label{subsec:sensing}

The sensing layer realizes the data-acquisition concerns of the SPR paradigm: selection of sensor types, number and locations \citep{farrar2013shm}. For the planned SHS braced-frame testbed, a heterogeneous, spatially distributed sensor network is foreseen, comprising electrical strain gauges (local strain and, through the constitutive law, stress), triaxial accelerometers (dynamic response for modal identification), temperature sensors (for thermal compensation and data normalisation), displacement sensors (global deflections), and inclinometers (global stability and rigid-body rotation). The planned channels and their monitoring roles are listed in Table~\ref{tab:sensors}. This heterogeneity is deliberate: data fusion across both homogeneous arrays (for example, a distributed strain-gauge array) and complementary modalities (acceleration combined with temperature and operational parameters) improves damage-detection fidelity and supports separation of benign variability from genuine damage \citep{farrar2013shm}. The number and placement of sensors are not chosen ad hoc but determined by preliminary numerical analysis of the finite-element baseline, which identifies the degrees of freedom most informative about the target modes and the most likely damage scenarios. This is consistent with sparse and optimal sensor-placement strategies (for example, QR- and DEIM-based selection and gappy POD) that allow reconstruction of the full response field from a limited number of measurements \citep{brunton2019data,tan2020sensorplacement,hassani2023sensorplacement}, and it respects the axiomatic principle that the length and time scales of the damage of interest dictate the required sensing-system properties \citep{farrar2013shm}. A fundamental constraint inherited at this layer is that sensors cannot measure damage directly; they measure response quantities from which damage information must subsequently be extracted \citep{farrar2013shm}.

\subsection{Data acquisition and signal-conditioning layer}\label{subsec:acquisition}

The acquisition layer converts the analog physical quantities sensed above into clean, synchronized and meaningful digital streams. A central DAQ system continuously samples the strain, acceleration, temperature, displacement and inclination channels at rates sufficient to resolve the highest natural frequencies of interest, in keeping with the snapshot-pair sampling requirements of data-driven modal methods \citep{brunton2019data}. Signal conditioning includes amplification, anti-alias filtering, bridge completion and excitation for the strain gauges, and time synchronization across channels so that mode shapes can be reconstructed from phase-consistent measurements. This layer also performs the data normalisation, cleansing, compression and fusion processes that are inherent to steps two through four of the SPR paradigm \citep{farrar2013shm}. Temperature channels are used here for thermal compensation, so that response changes caused by environmental variation, the canonical example being temperature-induced stiffness change, are not later misinterpreted as damage and reported as false positives \citep{farrar2013shm,figueiredo2010machine}. Temperature is the variability source this layer measures and removes directly; humidity, wind and randomly varying operational loads perturb the response as well, and because the testbed does not instrument them, they are left to the statistical normalization performed downstream in the cognition layer. Dimensionality reduction begins at this stage: the singular value decomposition and principal component analysis provide a numerically stable, data-driven means of compressing high-dimensional strain and acceleration streams into a few dominant coherent structures that define a low-dimensional baseline of healthy behavior \citep{brunton2019data}.

\subsection{Physics-based digital-twin layer}\label{subsec:digitaltwin}

At the core of the architecture is a finite-element model (FEM) that serves as the structure's physical baseline. The FEM idealizes the braced frame using bar and beam/frame elements appropriate to skeletal steel structures, with element stiffness matrices assembled into the global stiffness matrix and with explicit modeling decisions, notably the representation of bolted beam-column connections as semi-rigid rotational springs whose stiffness is an updatable parameter bounded by the pinned and rigid idealizations, treated as part of the mathematical model that the analysis actually solves \citep{cook2002fea}. The model supports the full suite of analyses required to establish baseline behavior: static analysis $[\mathbf{K}]\{\mathbf{D}\}=\{\mathbf{R}\}$, modal analysis via the undamped eigenproblem $([\mathbf{K}]-\omega^2[\mathbf{M}])\{\mathbf{D}\}=\{\mathbf{0}\}$ yielding natural frequencies and mode shapes, harmonic (frequency-response) analysis for the planned rotating-unbalance excitation, which the FE literature explicitly recognizes as a harmonic source applied by rotating machinery \citep{cook2002fea}, and transient response history. Mass is represented through consistent or lumped mass matrices and energy dissipation through Rayleigh proportional damping, $[\mathbf{C}]=a[\mathbf{M}]+b[\mathbf{K}]$ \citep{cook2002fea,chopra2017dynamics}. This model becomes a digital twin once it is continuously fed by sensor data and its predictions are compared in real time against measured responses \citep{chiachio2022structural,zou2026deep}. Because natural frequencies and mode shapes depend only on mass and stiffness \citep{chopra2017dynamics}, any stiffness loss, from progressive damage or bolted-connection loosening, manifests as a measurable shift in modal parameters and as a growing deviation between predicted and observed response. Model updating closes the gap: experimentally identified parameters are fed back to recalibrate the FEM, embodying the principle that finite-element analysis is simulation rather than reality and must be correlated with experiment \citep{cook2002fea}. Reduced-order representations, obtained through static condensation and substructuring \citep{cook2002fea} or POD-Galerkin projection \citep{brunton2019data}, make this comparison fast enough for real-time, streaming operation \citep{li2025edge}.

\subsection{Cognition layer}\label{subsec:cognition}

The cognition layer is where the structure interprets, rather than merely records, its own state, and it is organized along the Rytter damage-state hierarchy of existence, location, type, severity and prognosis \citep{farrar2013shm}. Perception is achieved through single-class models trained only on healthy data: autoencoders learn an unsupervised latent representation of normal behavior and flag abnormal conditions through growing reconstruction error \citep{brunton2019data}, while LSTM recurrent networks, used in a self-supervised forecasting mode, learn the healthy temporal evolution of the vibration and strain streams and flag deviations in sequential response \citep{brunton2019data,naser2023ml}. These novelty- and outlier-detection approaches answer the existence (and often location) questions in an unsupervised mode, consistent with the SHM axiom that existence and location can be found without labels whereas type and severity generally require supervision \citep{farrar2013shm,sohn2003statistical,wangcha2022unsupervised}. Diagnosis and damage-level classification are then handled by supervised learners, Random Forest, Support Vector Machines and Artificial Neural Networks, trained on labeled scenarios such as minor/major/collapse-style categories \citep{naser2023ml,parisi2022steel}. Because labeled damaged-state data are scarce when intentionally damaging a high-value structure is impractical \citep{farrar2013shm}, the validated FEM doubles as a data source: parametric simulation generates augmented, hybrid (real-plus-simulated) training sets \citep{naser2023ml}. The features consumed here are explicitly physics-derived: natural frequencies, mode shapes, damping ratios (estimable by half-power bandwidth, \citealp{chopra2017dynamics}), residuals between FEM prediction and measurement, and SVD/DMD modal descriptors \citep{brunton2019data}, so that physics and machine learning co-operate rather than compete \citep{avci2021review}.

\subsection{Decision and reporting layer}\label{subsec:decision}

The uppermost layer converts the cognition layer's quantified damage state into autonomous, actionable communication, realizing the operational-evaluation purpose of the SPR paradigm of life-safety and economic justification \citep{farrar2013shm}. Diagnostic and prognostic outputs are framed to minimize false-positive and false-negative diagnoses, whose costs may be weighted asymmetrically \citep{farrar2013shm}, and are made interpretable and auditable through explainable-ML techniques such as permutation feature importance, SHAP and partial-dependence plots that trace each alarm back to physically checkable indicators \citep{naser2023ml}. The structure thereby reports its own integrity, supporting the shift from time-based to condition-based and predictive maintenance \citep{farrar2013shm,mariniello2025enhancing}. Feedback flows downward, but under explicit governance, so that slow degradation is not silently absorbed into the notion of health. The reference healthy baseline is immutable and versioned; any adaptation of the working baseline, the FEM or the cognition models is conditioned on independent confirmation, monitored for concept drift, gated by human approval for safety-critical changes, and reversible through a rollback mechanism. Only validated, non-damaging changes, such as recalibration for a new operational or environmental regime, are propagated. This governed feedback closes the perceive--process--interpret--report loop that distinguishes a self-aware steel structure from conventional data-collecting SHM.
